% Source: https://artofproblemsolving.com/community/c3h1390147_slope_derivation
% Stable identity: aops-topic-1390147-q01
\textbf{(a)} Take two pairs of distinct points of $\ell$, say $P_1,P_2$ and $Q_1,Q_2$, and draw the slope triangles with horizontal and vertical legs.

\begin{asy}
usepackage("amsmath");
size(6cm);
defaultpen(fontsize(10pt));
pair P1=(1,2), P2=(4,6), Q1=(4.75,7), Q2=(6.25,9);
draw((0,0)--(7.5,0),Arrow); draw((0,0)--(0,10.6),Arrow);
label("$x$",(7.5,0),E); label("$y$",(0,10.6),N);
draw((0.25,1)--(7,10),linewidth(0.9));
draw(P1--(4,2)--P2,red);
draw((3.75,2)--(3.75,2.25)--(4,2.25));
draw(Q1--(6.25,7)--Q2,blue+dashed);
draw((6.05,7)--(6.05,7.2)--(6.25,7.2),blue);
dot(P1);dot(P2);dot(Q1);dot(Q2);
label("$P_1$",P1,NW); label("$P_2$",P2,NW);
label("$Q_1$",Q1,NW); label("$Q_2$",Q2,NW);
label("$x_2-x_1=3$",(2.5,2),S,red); label("$y_2-y_1=4$",(4,4),E,red);
label("\mbox{}",(min(currentpicture,user=true)+max(currentpicture,user=true))/2,black); // keeps the text after the figure black when LaTeX typesets the labels inline
\end{asy}

If $\ell$ is horizontal, every ratio is $0$. Otherwise the two right triangles have a horizontal leg, a vertical leg, and a hypotenuse along $\ell$, so the angle between the hypotenuse and the horizontal leg is the same in both (corresponding angles of the parallel horizontal legs cut by $\ell$). By AA the triangles are similar, so the ratios of vertical leg to horizontal leg are equal:
\[
\dfrac{|y_2-y_1|}{|x_2-x_1|}=\dfrac{|y_2'-y_1'|}{|x_2'-x_1'|}.
\]
The signs also agree: along a line, $y$ increases with $x$ throughout (if the line rises) or decreases throughout (if it falls), so $y_2-y_1$ and $x_2-x_1$ have the same sign or opposite signs for every pair. (Swapping $P_1$ and $P_2$ changes both differences in sign and leaves the ratio unchanged.) Hence the ratio is a property of the line alone. The hypotenuse $\sqrt{(x_2-x_1)^2+(y_2-y_1)^2}$ is the distance formula; it changes when the points move along $\ell$ (for example, doubling the run doubles it), so it cannot measure steepness.

\textbf{(b)} The slope is $\dfrac{6-2}{4-1}=\boxed{\dfrac{4}{3}}$, and the distance is $\sqrt{3^2+4^2}=\boxed{5}$.
